\section{总结}

在本文中，我们证明了 quasi-PDFs 可以看作 PDFs 与部分子在强子静止系中
原始动量分布的混合体。
在此图像下，部分子的 $k_3$ 动量
既来自强子整体的
运动，也来自静止系中的原始
动量分布。quasi-PDFs 复杂的卷积本质
要求使用 $p_3 \gtrsim $ \mbox{3 GeV} 的动量才能消除
原始动量分布的效应并
足够接近 PDF 极限。

作为替代方案，我们建议使用
pseudo-PDFs ${\cal P} (x, z_3^2)  $，它把光前 PDFs 推广到类空间隔。
通过 Fourier 变换，它们与 Ioffe 时间分布
${\cal M} (\nu, z_3^2) $ 相联系，后者由写成 $\nu = p_3 z_3$ 和 $z_3^2$
函数的一般矩阵元素给出。
pseudo-PDFs 的优点在于：第一，它们具有与 PDFs 相同的
$-1 \leq x \leq 1$ 支撑；第二，在小 $z_3^2$ 处其
$z_3^2$ 依赖性由通常的演化方程支配。

构成 Ioffe 时间分布的比值
${\cal M} (\nu, z_3^2)/{\cal M} (0, z_3^2)$，
可以把原始静止系分布产生的绝大部分 $z_3^2$ 依赖性约去。
此外，取此比值还能排除来自 $\hat E (0,z_3;A)$ 链接格点重正化的
\mbox{$z_3^2$ 依赖} 因子——该因子给 quasi-PDFs 的
格点计算带来困难（见例如 \cite{Ishikawa:2016znu}）。

检验用 pseudo-PDFs 从格点提取 PDFs 的效率，
是对未来研究的挑战。

事实上，就在本文处于评审过程期间，
一篇基于本文思想的实际格点计算 \cite{Orginos:2017kos} 已经完成。
它清楚地证明了静止系函数 ${\cal M} (0,z_3^2)$ 的 $z_3$ 依赖性中
存在线性成分，可以归因于规范链接产生的
$Z(z_3^2/a^2) \sim e^{-c|z_3|/a}$ 行为。研究还观察到，
比值 ${\cal M} (Pz_3 , z_3^2)/{\cal M} (0, z_3^2)$
对 $z_3$ 呈 Gaussian 型行为，这表明
进入 ${\mathfrak M} (Pz_3, z_3^2)$ 比值分子和分母的
$Z(z_3^2/a^2)$ 因子如预期般相消了。

此外还发现，当以 $\nu$ 和 $z_3$ 为变量作图时，
约化分布 ${\mathfrak M} (\nu, z_3^2)$ 的数据对 $z_3^2$ 的依赖性非常温和。
这一观察表明，
${\cal M} (\nu, z_3^2)$ 的 \mbox{$z_3^2$ 依赖性} 中的软部分
已被静止系密度 ${\cal M} (0, z_3^2)$ 抵消。
该现象对应于 TMD ${\cal F} (x, k_\perp^2)$ 软部分的
$x$ 依赖性与 $k_\perp$ 依赖性的因子化。

研究还证明，在小 $z_3 \leq 4 \,a$ 处残余的 \mbox{$z_3$ 依赖性}
可以用微扰演化解释，相应的 $\alpha_s$ 值
对应于 $\alpha_s/\pi =0.1$。
